% GENERATED FILE. Edit canonical lessons, then run npm run generate:books.
% canonical-source-hash: fnv1a64:43831bfa253a05c9
% canonical-lessons: PA-C338-ijazat, PA-C338-manha, PA-C338-manzur, PA-C338-hukam, PA-C338-tokna

\chapter{Is It Allowed?}
\label{ch:pa-a2-is-it-allowed}
\hlchaptermodality{338}

\begin{chapteropening}
\textbf{By the end of this chapter:} \emph{I can ask for permission and say what is not allowed.}

Complete the last lesson of chapter 338: I can ask for permission and say what is not allowed.
\end{chapteropening}

\section[\texorpdfstring{\pa{ਇਜਾਜ਼ਤ} (ijāzat)}{ਇਜਾਜ਼ਤ (ijāzat)}]{\pa{ਇਜਾਜ਼ਤ} (ijāzat) — permission}

\label{lesson:PA-C338-ijazat}



\begin{quote}
Before the new one: say the Punjabi for a bit (to soften a request), then the Punjabi for to send for.
\end{quote}

\subsection*{You'll want to know: \pa{ਇਜਾਜ਼ਤ}}
\textbf{\pa{ਇਜਾਜ਼ਤ}} — \emph{ijāzat} — \textquotedblleft{}permission\textquotedblright{}.

\begin{quote}
\pa{ਅਧਿਆਪਕ} \pa{ਇਜਾਜ਼ਤ} \pa{ਨਹੀਂ} \pa{ਦਿੰਦਾ} \pa{ਹੈ।}
\end{quote}

\begin{quote}
\emph{adhiāpak ijāzat nahīṁ dindā hai.}
\end{quote}

\begin{quote}
The teacher does not give permission.
\end{quote}

A feminine noun.

\begin{cousinweb}
\textbf{\pa{ਇਜਾਜ਼ਤ}} (\emph{ijāzat}) is a loan, usually traced through Persian to Arabic \emph{ijāza}, \textquotedblleft{}permission\textquotedblright{}.
\end{cousinweb}

\begin{grammarlens}[title={What you've built}]
One more word for asking and answering.
\end{grammarlens}

\subsection*{Your turn}
\begin{itemize}
\raggedright
  \item \emph{Say it:} \emph{ijāzat}
  \item \emph{Say it:} \emph{ijāzat}, once more
  \item \emph{Say it:} \emph{adhiāpak ijāzat nahīṁ dindā hai.} — The teacher does not give permission.
  \item \emph{Recall:} say the Punjabi for a bit (to soften a request), then the Punjabi for to send for, then say \emph{ijāzat} again
\end{itemize}

\begin{tcolorbox}[breakable,colback=teal!4,colframe=teal!35!black,title={Before you move on}]
What does \pa{ਇਜਾਜ਼ਤ} mean? (\textquotedblleft{}Permission\textquotedblright{}.) Say it once more.
\end{tcolorbox}
\section[\texorpdfstring{\pa{ਮਨ੍ਹਾ} (manhā)}{ਮਨ੍ਹਾ (manhā)}]{\pa{ਮਨ੍ਹਾ} (manhā) — forbidden; not allowed}

\label{lesson:PA-C338-manha}



\begin{quote}
Before the new one: say the Punjabi for to send for, then the Punjabi for permission.
\end{quote}

\subsection*{You'll want to know: \pa{ਮਨ੍ਹਾ}}
\textbf{\pa{ਮਨ੍ਹਾ}} — \emph{manhā} — \textquotedblleft{}forbidden; not allowed\textquotedblright{}.

\begin{quote}
\pa{ਕੀ} \pa{ਬਾਗ਼} \pa{ਵਿੱਚ} \pa{ਖੇਡਣਾ} \pa{ਮਨ੍ਹਾ} \pa{ਹੈ}?
\end{quote}

\begin{quote}
\emph{kī bāgh vich kheḍṇā manhā hai?}
\end{quote}

\begin{quote}
Is playing in the park forbidden?
\end{quote}

\begin{cousinweb}
\textbf{\pa{ਮਨ੍ਹਾ}} (\emph{manhā}) is a loan, usually traced through Persian to Arabic \emph{manʿ}, \textquotedblleft{}prevention\textquotedblright{}.
\end{cousinweb}

\begin{grammarlens}[title={What you've built}]
One more word for asking and answering.
\end{grammarlens}

\subsection*{Your turn}
\begin{itemize}
\raggedright
  \item \emph{Say it:} \emph{manhā}
  \item \emph{Say it:} \emph{manhā}, once more
  \item \emph{Say it:} \emph{kī bāgh vich kheḍṇā manhā hai?} — Is playing in the park forbidden?
  \item \emph{Recall:} say the Punjabi for to send for, then the Punjabi for permission, then say \emph{manhā} again
  \item \emph{Write it:} copy \textbf{\pa{ਮਨ੍ਹਾ}} once, checking the whole word
\end{itemize}

\begin{tcolorbox}[breakable,colback=teal!4,colframe=teal!35!black,title={Before you move on}]
What does \pa{ਮਨ੍ਹਾ} mean? (\textquotedblleft{}Forbidden; not allowed\textquotedblright{}.) Say it once more.
\end{tcolorbox}
\section[\texorpdfstring{\pa{ਮਨਜ਼ੂਰ} (manzūr)}{ਮਨਜ਼ੂਰ (manzūr)}]{\pa{ਮਨਜ਼ੂਰ} (manzūr) — approved; accepted}

\label{lesson:PA-C338-manzur}



\begin{quote}
Before the new one: say the Punjabi for permission, then the Punjabi for forbidden.
\end{quote}

\subsection*{You'll want to know: \pa{ਮਨਜ਼ੂਰ}}
\textbf{\pa{ਮਨਜ਼ੂਰ}} — \emph{manzūr} — \textquotedblleft{}approved; accepted\textquotedblright{}.

\begin{quote}
\pa{ਕੀ} \pa{ਮੇਰੀ} \pa{ਅਰਜ਼ੀ} \pa{ਮਨਜ਼ੂਰ} \pa{ਹੈ}?
\end{quote}

\begin{quote}
\emph{kī merī arzī manzūr hai?}
\end{quote}

\begin{quote}
Has my application been approved?
\end{quote}

\begin{cousinweb}
\textbf{\pa{ਮਨਜ਼ੂਰ}} (\emph{manzūr}) is a loan, usually traced through Persian to Arabic \emph{manẓūr}, \textquotedblleft{}looked at\textquotedblright{}.
\end{cousinweb}

\begin{grammarlens}[title={What you've built}]
One more word for asking and answering.
\end{grammarlens}

\subsection*{Your turn}
\begin{itemize}
\raggedright
  \item \emph{Say it:} \emph{manzūr}
  \item \emph{Say it:} \emph{manzūr}, once more
  \item \emph{Say it:} \emph{kī merī arzī manzūr hai?} — Has my application been approved?
  \item \emph{Recall:} say the Punjabi for permission, then the Punjabi for forbidden, then say \emph{manzūr} again
  \item \emph{Write it:} copy \textbf{\pa{ਮਨਜ਼ੂਰ}} once, checking the whole word
\end{itemize}

\begin{tcolorbox}[breakable,colback=teal!4,colframe=teal!35!black,title={Before you move on}]
What does \pa{ਮਨਜ਼ੂਰ} mean? (\textquotedblleft{}Approved; accepted\textquotedblright{}.) Say it once more.
\end{tcolorbox}
\section[\texorpdfstring{\pa{ਹੁਕਮ} (hukam)}{ਹੁਕਮ (hukam)}]{\pa{ਹੁਕਮ} (hukam) — an order; a command}

\label{lesson:PA-C338-hukam}



\begin{quote}
Before the new one: say the Punjabi for forbidden, then the Punjabi for approved.
\end{quote}

\subsection*{You'll want to know: \pa{ਹੁਕਮ}}
\textbf{\pa{ਹੁਕਮ}} — \emph{hukam} — \textquotedblleft{}an order; a command\textquotedblright{}.

\begin{quote}
\pa{ਮੁੰਡਾ} \pa{ਪਿਤਾ} \pa{ਦਾ} \pa{ਹੁਕਮ} \pa{ਨਹੀਂ} \pa{ਮੰਨਦਾ} \pa{ਹੈ।}
\end{quote}

\begin{quote}
\emph{muṇḍā pitā dā hukam nahīṁ manndā hai.}
\end{quote}

\begin{quote}
The boy does not obey his father's order.
\end{quote}

A masculine noun.

\begin{cousinweb}
\textbf{\pa{ਹੁਕਮ}} (\emph{hukam}) is a loan, usually traced through Persian to Arabic \emph{ḥukm}, \textquotedblleft{}judgement\textquotedblright{}.
\end{cousinweb}

\begin{grammarlens}[title={What you've built}]
One more word for asking and answering.
\end{grammarlens}

\subsection*{Your turn}
\begin{itemize}
\raggedright
  \item \emph{Say it:} \emph{hukam}
  \item \emph{Say it:} \emph{hukam}, once more
  \item \emph{Say it:} \emph{muṇḍā pitā dā hukam nahīṁ manndā hai.} — The boy does not obey his father's order.
  \item \emph{Recall:} say the Punjabi for forbidden, then the Punjabi for approved, then say \emph{hukam} again
  \item \emph{Write it:} copy \textbf{\pa{ਹੁਕਮ}} once, checking the whole word
\end{itemize}

\begin{tcolorbox}[breakable,colback=teal!4,colframe=teal!35!black,title={Before you move on}]
What does \pa{ਹੁਕਮ} mean? (\textquotedblleft{}An order; a command\textquotedblright{}.) Say it once more.
\end{tcolorbox}
\section[\texorpdfstring{\pa{ਟੋਕਣਾ} (ṭokṇā)}{ਟੋਕਣਾ (ṭokṇā)}]{\pa{ਟੋਕਣਾ} (ṭokṇā) — to interrupt; to stop someone with a remark}

\label{lesson:PA-C338-tokna}



\begin{quote}
Before the new one: say the Punjabi for approved, then the Punjabi for an order.
\end{quote}

\subsection*{You'll want to know: \pa{ਟੋਕਣਾ}}
\textbf{\pa{ਟੋਕਣਾ}} — \emph{ṭokṇā} — \textquotedblleft{}to interrupt; to stop someone with a remark\textquotedblright{}.

\begin{quote}
\pa{ਮਿਹਰਬਾਨੀ} \pa{ਕਰਕੇ} \pa{ਮੈਨੂੰ} \pa{ਨਾ} \pa{ਟੋਕੋ।}
\end{quote}

\begin{quote}
\emph{miharbānī karke mainū̃ nā ṭoko.}
\end{quote}

\begin{quote}
Please don't interrupt me.
\end{quote}

Past (I …): \pa{ਮੈਂ} \pa{ਅਮਨ} \pa{ਨੂੰ} \pa{ਟੋਕਿਆ}, \emph{maiṁ Aman nū̃ ṭokiā}.

\begin{cousinweb}
\textbf{\pa{ਟੋਕਣਾ}} (\emph{ṭokṇā}) is inherited from a reconstructed \emph{*ṭōkk}, \textquotedblleft{}strike, bite\textquotedblright{}; CDIAL lists Punjabi \emph{ṭokṇā}, \textquotedblleft{}to hinder\textquotedblright{}, under it (CDIAL 5476).
\end{cousinweb}

\begin{grammarlens}[title={What you've built}]
One more word for asking and answering.
\end{grammarlens}

\subsection*{Your turn}
\begin{itemize}
\raggedright
  \item \emph{Say it:} \emph{ṭokṇā}
  \item \emph{Say it:} \emph{ṭokṇā}, once more
  \item \emph{Say it:} \emph{miharbānī karke mainū̃ nā ṭoko.} — Please don't interrupt me.
  \item \emph{Recall:} say the Punjabi for approved, then the Punjabi for an order, then say \emph{ṭokṇā} again
  \item \emph{Write it:} copy \textbf{\pa{ਟੋਕਣਾ}} once, checking the whole word
\end{itemize}

\begin{tcolorbox}[breakable,colback=teal!4,colframe=teal!35!black,title={Before you move on}]
What does \pa{ਟੋਕਣਾ} mean? (\textquotedblleft{}To interrupt; to stop someone with a remark\textquotedblright{}.) Say it once more.
\end{tcolorbox}
